\section{Estructura del firmware}
\label{sec:estructura_firmware}

\subsection{Código fuente}
\label{sec:codigo_fuente}

El código completo del proyecto está disponible en un repositorio público de GitHub:

\begin{center}
\url{https://github.com/sevastianamoroso/proyecto-mano}
\end{center}

El repositorio contiene el firmware del guante (\texttt{Transmisor/}) y de la mano (\texttt{Receptor/}), el receptor modificado para la medición de latencia (\texttt{Receptor\_Latencia/}), el plan y el registro de pruebas, la documentación del firmware y las fuentes de este informe. Cada versión del firmware descrita en la sección~\ref{sec:evolucion} está marcada con una etiqueta de git (\texttt{v1.0.0} a \texttt{v2.3.0}), de modo que cualquiera de ellas puede recuperarse y volver a grabarse en los ESP32. Este apéndice describe la versión 2.3.0.

\subsection{Cómo leer los diagramas}

El firmware está escrito en C++ para Arduino, pero \textbf{no define clases propias}: está organizado en variables globales, funciones y tareas de FreeRTOS, repartidos en dos programas (\texttt{Transmisor.ino} y \texttt{Receptor.ino}) que comparten el archivo \texttt{ProtocoloBrazo.h}. Para mostrar su estructura se usa la notación de los diagramas de clases de UML, agrupando las variables y funciones que cumplen una misma responsabilidad en un \emph{módulo}. Cada caja tiene tres partes:

\begin{itemize}
  \item \textbf{Encabezado:} entre comillas angulares, el tipo de elemento (\emph{estereotipo}): tarea de FreeRTOS, \emph{callback}, estructura, enumeración, etc. Debajo va el nombre del módulo.
  \item \textbf{Atributos:} las variables y constantes del módulo, con su tipo.
  \item \textbf{Operaciones:} las funciones que lo forman.
\end{itemize}

Los nombres de los módulos que no existen como tales en el código (por ejemplo, \emph{EstimadorMuñeca}) son agrupaciones para la explicación; los de variables y funciones son los reales. Las flechas tienen el siguiente significado:

\begin{center}
\begin{tikzpicture}[every node/.style={font=\small}]
  \draw[asoc] (0,0) -- (1.5,0) node[right] {asociación: crea, llama o envía a};
  \draw[usa]  (0,-0.6) -- (1.5,-0.6) node[right] {dependencia: usa o lee};
  \draw[comp] (0,-1.2) -- (1.5,-1.2) node[right] {composición: es parte de (el rombo va del lado del contenedor)};
\end{tikzpicture}
\end{center}

% =====================================================================
\subsection{Vista general}

La figura~\ref{fig:fw_general} muestra los dos programas y lo que comparten. El transmisor (guante) mide los sensores, calcula los ángulos de los servos y los envía por ESP-NOW 50 veces por segundo. El receptor (mano) valida cada paquete, lo pasa a su tarea de control por una cola de FreeRTOS y mueve los servos de forma suave. Ambos programas usan la misma definición del paquete, \texttt{MensajeBrazo}, por lo que el formato no puede quedar desincronizado entre ellos.

\begin{figure}[H]
\centering
\begin{tikzpicture}[node distance=5mm]
  % Transmisor
  \node[modulo=ctx] (tmain) {Programa principal\\\scriptsize\texttt{setup()}, \texttt{loop()}};
  \node[modulo=crtos, below=of tmain] (ttask) {TareaTransmision\\\scriptsize 50~Hz};
  \node[modulo=ctx, below=of ttask] (timu) {EstimadorMuñeca\\\scriptsize MPU6050};
  \node[modulo=ctx, below=of timu] (tdedos) {LecturaDedos\\\scriptsize potenciómetros};
  \node[modulo=ctx, below=of tdedos] (tenl) {EnlaceTx\\\scriptsize ESP-NOW};
  \begin{scope}[on background layer]
    \node[paquete, fill=ctx!30, fit=(tmain)(tenl), inner xsep=16pt, label={[font=\bfseries]above:Transmisor.ino (guante)}] (ptx) {};
  \end{scope}

  % Receptor
  \node[modulo=crx, right=6.2cm of tmain] (rmain) {Programa principal\\\scriptsize\texttt{setup()}, \texttt{loop()}};
  \node[modulo=crx, below=of rmain] (rcb) {OnDataRecv\\\scriptsize \emph{callback} de recepción};
  \node[modulo=crtos, below=of rcb] (rcola) {colaMensajesBrazo\\\scriptsize cola de 1 elemento};
  \node[modulo=crtos, below=of rcola] (rtask) {TareaControlActuadores\\\scriptsize máquina de estados};
  \node[modulo=crx, below=of rtask] (ract) {Actuadores\\\scriptsize 7 servos};
  \begin{scope}[on background layer]
    \node[paquete, fill=crx!30, fit=(rmain)(ract), inner xsep=16pt, label={[font=\bfseries]above:Receptor.ino (mano)}] (prx) {};
  \end{scope}

  % Compartido
  \node[modulo=cshared, below=1.4cm of $(ptx.south)!0.5!(prx.south)$] (msg) {MensajeBrazo\\\scriptsize paquete de 19 bytes};
  \begin{scope}[on background layer]
    \node[paquete, fill=cshared!40, fit=(msg), label={[font=\bfseries]below:ProtocoloBrazo.h (compartido)}] {};
  \end{scope}

  % Transmisor: la tarea usa los tres módulos (bus por la izquierda, dentro del paquete)
  \draw[asoc] (tmain) -- (ttask);
  \draw[usa] (ttask.west) -- ++(-0.35,0) |- (timu.west);
  \draw[usa] ([xshift=-0.35cm]ttask.west) |- (tdedos.west);
  \draw[usa] ([xshift=-0.35cm]ttask.west) |- (tenl.west);
  % Enlace de radio
  \draw[asoc, line width=1.4pt] (tenl.east) -- node[etq, above, sloped] {ESP-NOW (radio)} (rcb.west);
  % Receptor: cadena de recepción y creación de la tarea (por la derecha, dentro del paquete)
  \draw[asoc] (rmain) -- (rcb);
  \draw[asoc] (rcb) -- (rcola);
  \draw[asoc] (rcola) -- (rtask);
  \draw[asoc] (rtask) -- (ract);
  \draw[asoc] (rmain.east) -- ++(0.35,0) |- (rtask.east);
  % Ambos programas incluyen el protocolo
  \draw[usa] (ptx.south) |- node[etq, pos=0.75, above] {incluye} (msg.west);
  \draw[usa] (prx.south) |- node[etq, pos=0.75, above] {incluye} (msg.east);
\end{tikzpicture}
\caption{Vista general del firmware. En violeta, los elementos de FreeRTOS (tareas y cola).}
\label{fig:fw_general}
\end{figure}

\begin{table}[H]
\centering
\small
\begin{tabularx}{\textwidth}{l l X}
\toprule
\textbf{Módulo} & \textbf{Programa} & \textbf{Responsabilidad} \\
\midrule
Programa principal & Tx & Configura el hardware, la radio y la tarea; en \texttt{loop()} atiende el comando de calibración. \\
TareaTransmision & Tx & Ciclo de \SI{20}{ms}: lee la IMU y los dedos, calcula los ángulos, arma el paquete y lo envía. \\
EstimadorMuñeca & Tx & Lectura del MPU6050 y filtro complementario para \emph{pitch} y \emph{roll}. \\
LecturaDedos & Tx & Lectura de los potenciómetros con sobremuestreo y calibración de dos puntos guardada en memoria no volátil. \\
EnlaceTx & Tx & Envío por ESP-NOW y registro del resultado del último envío. \\
MensajeBrazo & ambos & Formato del paquete, cálculo y validación del \emph{checksum}. \\
OnDataRecv & Rx & Descarta paquetes inválidos y deja el válido en la cola. \\
TareaControlActuadores & Rx & Máquina de estados, \emph{failsafe}, limitador de velocidad y banda muerta. \\
Actuadores & Rx & Manejo de los siete servos: acople, desacople e inversión de sentido. \\
\bottomrule
\end{tabularx}
\caption{Responsabilidad de cada módulo.}
\label{tab:fw_modulos}
\end{table}

% =====================================================================
\newpage
\subsection{Transmisor}

La tarea \texttt{tareaTransmision} es el centro del transmisor. Se ejecuta cada \SI{20}{ms} en el núcleo~1 y en cada ciclo usa, en este orden, al estimador de la muñeca, a la lectura de los dedos y al enlace. El filtro complementario y la conversión de ángulos de la mano a ángulos de servo están escritos dentro del ciclo de la tarea; en el diagrama se los ubica en \emph{EstimadorMuñeca} porque operan sobre sus variables.

La calibración de los dedos corre en \texttt{loop()}, fuera de la tarea. Para que ambas no lean el ADC al mismo tiempo, la calibración activa la bandera \texttt{calibrando}; mientras está activa, la tarea sigue actualizando el filtro de la muñeca pero no lee los dedos, no transmite y no imprime telemetría.

\begin{figure}[H]
\centering
\resizebox{\textwidth}{!}{%
\begin{tikzpicture}
  \clase[clase=ctx]{main}{programa principal}{Transmisor.ino}
    {—}
    {setup()\\loop() \textrm{: comando} 'c'}
  \clase[clase=crtos, at={(6.4,0)}]{task}{tarea · núcleo 1 · prioridad 2}{TareaTransmision}
    {PERIODO\_CONTROL\_MS = 20\\
     xLastWakeTime : TickType\_t\\
     ticksTelemetria : uint32\_t\\
     paqueteSalida : MensajeBrazo\\
     contadorSecuencia : uint8\_t}
    {tareaTransmision()\\
     \textrm{\itshape ciclo: IMU $\rightarrow$ filtro $\rightarrow$ mapeo $\rightarrow$ dedos}\\
     \textrm{\itshape $\rightarrow$ paquete $\rightarrow$ envío $\rightarrow$ telemetría (5 Hz)}}
  \clase[clase=cshared, at={(12.8,0)}]{msg}{struct · ProtocoloBrazo.h · 19 bytes}{MensajeBrazo}
    {magicHeader : uint16\_t = 0xAA55\\
     secuencia : uint8\_t\\
     anguloDedos[5] : uint16\_t\\
     muniecaVertical : int16\_t\\
     muniecaRotacional : int16\_t\\
     checksum : uint16\_t}
    {calcularChecksumBrazo(msg) : uint16\_t\\
     validarPaqueteBrazo(msg) : bool}

  \clase[clase=ctx, at={(6.4,-5.4)}]{imu}{módulo}{EstimadorMuñeca}
    {mpu : Adafruit\_MPU6050\\
     anguloPitch, anguloRoll : float\\
     gyroBiasX, gyroBiasY : float\\
     tiempoAnteriorMicros : uint32\_t\\
     ALPHA\_COMPLEMENTARIO = 0.98\\
     TRIM\_PITCH\_DEG, TRIM\_ROLL\_DEG\\
     PITCH\_MANO\_*, SERVO\_V\_*, ROLL\_MANO\_MAX}
    {rutinaRecuperacionI2C(sda, scl)\\
     calibrarIMU()\\
     pitchAcc(a) : float\\
     rollAcc(a) : float}
  \clase[clase=ctx, at={(0,-5.4)}]{dedos}{módulo}{LecturaDedos}
    {PINES\_DEDOS[5] : int\\
     calMinADC[5], calMaxADC[5] : uint16\_t\\
     memoriaNVS : Preferences\\
     calibrando : volatile bool}
    {leerADCSobremuestreado(pin, muestras)\\
     ejecutarCalibracionInteractiva()\\
     cargarCalibracionNVS()\\
     guardarCalibracionNVS()}
  \clase[clase=ctx, at={(12.8,-5.4)}]{enl}{módulo · ESP-NOW}{EnlaceTx}
    {macReceptor[6] : uint8\_t\\
     infoReceptor : esp\_now\_peer\_info\_t\\
     ultimoEnvioExitoso : volatile bool}
    {OnDataSent(info, status) \textrm{«callback»}\\
     esp\_now\_send(...) \textrm{(desde la tarea)}}

  \draw[asoc] (main.east |- task.text split) -- node[etq, above] {crea} (task.west |- task.text split);
  \draw[asoc] (task.east |- msg.text split) -- node[etq, above] {arma} (msg.west |- msg.text split);
  % La tarea usa los tres módulos; LecturaDedos queda debajo del programa principal para que
  % la flecha del comando 'c' baje recta sin cruzar el bus.
  \coordinate (bus) at (6.4,-4.4);
  \draw[usa] (task.south) -- node[etq, pos=0.8, right] {usa} (imu.north);
  \draw[usa] (bus) -| ([xshift=1.2cm]dedos.north);
  \draw[usa] (bus) -| ([xshift=-0.6cm]enl.north);
  \draw[asoc] ([xshift=-1.2cm]main.south) -- node[etq, pos=0.6, right] {comando 'c'} ([xshift=-1.2cm]dedos.north);
  \draw[usa] ([xshift=1.6cm]enl.north) -- node[etq, right] {envía} ([xshift=1.6cm]msg.south);
\end{tikzpicture}}
\caption{Estructura del transmisor.}
\label{fig:fw_tx}
\end{figure}

% =====================================================================
\subsection{Receptor}

En el receptor, el \emph{callback} \texttt{OnDataRecv} corre en el núcleo~0, dentro de la pila de Wi-Fi. Sólo verifica el tamaño y el \emph{checksum} del paquete y lo copia en \texttt{colaMensajesBrazo}, una cola de un único elemento que se sobrescribe: si llegan dos paquetes antes de que la tarea los procese, se conserva el más reciente. La tarea \texttt{tareaControlActuadores}, en el núcleo~1, espera en esa cola con un tiempo máximo de \SI{20}{ms}; así despierta en cuanto llega un paquete y, si no llega ninguno, igual sigue ejecutándose para supervisar el \emph{failsafe} y completar la trayectoria hacia la posición de reposo. Como sólo la tarea modifica las posiciones de los servos, no hay accesos simultáneos desde el \emph{callback}.

\begin{figure}[H]
\centering
\resizebox{\textwidth}{!}{%
\begin{tikzpicture}
  \clase[clase=crx]{cb}{callback ESP-NOW · núcleo 0}{OnDataRecv}
    {—}
    {OnDataRecv(info, datos, len)\\
     \textrm{\itshape 1. descarta si len $\neq$ 19}\\
     \textrm{\itshape 2. validarPaqueteBrazo()}\\
     \textrm{\itshape 3. xQueueOverwriteFromISR()}}
  \clase[clase=crtos, at={(6.4,0)}]{cola}{cola FreeRTOS}{colaMensajesBrazo}
    {longitud = 1\\
     elemento : MensajeBrazo\\
     \textrm{sobrescribe el anterior}}
    {xQueueOverwriteFromISR() \textrm{(escribe)}\\
     xQueueReceive() \textrm{(lee)}}
  \clase[clase=crtos, at={(12.8,0)}]{task}{tarea · núcleo 1 · prioridad 3}{TareaControlActuadores}
    {estadoActual : EstadoSistema\\
     marcaTiempoUltimoPaquete : uint32\_t\\
     tiempoEnHomeMs : uint32\_t\\
     ticksTelemetria : uint32\_t}
    {tareaControlActuadores()\\
     \textrm{\itshape 1. esperar paquete (máx. 20 ms)}\\
     \textrm{\itshape 2. supervisar failsafe (250 ms)}\\
     \textrm{\itshape 3. trayectoria: velocidad + banda muerta}\\
     \textrm{\itshape 4. desacople en reposo (2 s)}\\
     \textrm{\itshape 5. telemetría (5 Hz)}}

  \clase[clase=crx, at={(0,-3.4)}]{main}{programa principal}{Receptor.ino}
    {—}
    {setup()\\
     \textrm{\itshape posición Home, temporizadores PWM,}\\
     \textrm{\itshape cola, ESP-NOW y tarea}\\
     loop() \textrm{\itshape (vacío)}}
  \clase[clase=cshared, at={(6.4,-3.3)}]{est}{enumeración}{EstadoSistema}
    {ESTADO\_ESPERANDO\_SYNC\\
     ESTADO\_OPERACION\_NORMAL\\
     ESTADO\_FAILSAFE\_SIN\_SENAL}
    {—}
  \clase[clase=crx, at={(12.8,-5.6)}]{act}{módulo}{Actuadores}
    {servoDedos[5] : Servo \textrm{(SG90)}\\
     servoMuniecaVert, servoMuniecaRot : Servo\\
     posActualDedos[5], posObjetivoDedos[5]\\
     posActualMunieca*, posObjetivoMunieca*\\
     servosAcoplados : bool}
    {acoplarServos()\\
     desacoplarServos()\\
     anguloServoDedo(i, pos) : int}
  \clase[clase=cshared, at={(6.4,-6.1)}]{par}{constantes}{ParámetrosControl}
    {TIMEOUT\_FAILSAFE\_MS = 250\\
     DEADBAND\_DEDOS = 1.0, DEADBAND\_MUNIECA = 1.2\\
     VEL\_MAX\_DEDOS = 8.0 \textrm{°/ciclo}\\
     VEL\_MAX\_MUNIECA\_VERT = 4.0, \_ROT = 5.0\\
     HOME\_*, LIM\_*\\
     INVERTIR\_DEDO[5] : bool}
    {—}

  \draw[asoc] (cb.east |- cola.text split) -- node[etq, above] {escribe} (cola.west |- cola.text split);
  \draw[asoc] (cola.east |- cola.text split) -- node[etq, above] {lee} (task.west |- cola.text split);
  \draw[comp] (task.west |- est.east) -- node[etq, above] {estado} (est.east);
  \draw[asoc] ([xshift=1cm]task.south) -- node[etq, right] {mueve} ([xshift=1cm]act.north);
  \draw[usa] (act.west |- par.east) -- node[etq, above] {usa} (par.east);
  \draw[asoc] (main.north) -- node[etq, left] {registra} (cb.south);
\end{tikzpicture}}
\caption{Estructura del receptor.}
\label{fig:fw_rx}
\end{figure}

Las bibliotecas externas no se dibujan para no recargar los diagramas: \texttt{Adafruit\_MPU6050} (IMU), \texttt{Preferences} (memoria no volátil), \texttt{ESP32Servo} (PWM de los servos), \texttt{esp\_now} y \texttt{WiFi} (radio) y FreeRTOS (tareas y colas), incluida en el núcleo de Arduino para ESP32.
